\subsection{One coherent replacement over every Fermat level and the full metabelian quotient}

\begin{statement}{D.47}{Proved}{The closed relative-prime source, all Fermat homology maps, and its nonarithmetical metabelian action}
All subgroups below are closed in \(\widehat B\), and every
commutator subgroup denotes the closed subgroup generated
by commutators. Let
\[
 \begin{aligned}
 F_m&=\ker(\widehat F\to(\mathbb Z/m)^2),&
 L_m&=\widehat N\cap F_m,\\
 M&=\widehat N\cap\widehat F'
          =\bigcap_{m\ge1}L_m,&
 \mathcal H_\infty&=\widehat B/M .
 \end{aligned}                                                   \tag{D.47.1}
\]
The maps use the original ordered loop abelianization.
For every prime \(p\), write \(M^{(p)}\) for the maximal
continuous pro-\(p\) quotient of this compact group, and set
\[
 \begin{aligned}
 R_{\infty,p}&=\ker(M\to M^{(p)}),&
 \mathcal P_{\infty,p}&=\widehat B/R_{\infty,p},\\
 R_\infty&=\bigcap_pR_{\infty,p},&
 \mathcal V&=\widehat B/R_\infty .
 \end{aligned}                                                   \tag{D.47.2}
\]
There are exact marked identifications
\[
 \begin{aligned}
 \mathcal P_{\infty,p}
   &=\varprojlim_m
          \widehat B/\ker(L_m\to L_m^{(p)}),\\
 \mathcal V&=\prod_{\mathcal H_\infty}
                                     \mathcal P_{\infty,p},
       \qquad
 \ker(\mathcal V\to\mathcal H_\infty)
                         =\prod_pM^{(p)} .
 \end{aligned}                                                   \tag{D.47.3}
\]
The first inverse limit is over divisibility, with all maps
induced by the original \(\widehat B\). The maximal common
marked quotient of two different prime factors, or of one
factor and the joint remaining prime factors, is exactly
\(\mathcal H_\infty\). Both kernel inclusions in the first
line, and the actual simultaneous image in the second,
are proved below.

The base and all its finite levels are explicit. Put
\(H=q(\widehat F)\subset Q\), and \(J=H'\).
Then \(J\simeq Q_8\), \(|J|=8\), and
\[
 \mathcal H_\infty
   =(\widehat B/\widehat F')
                         \times_{Q/J}Q,\qquad |Q/J|=72.
                                                               \tag{D.47.4}
\]
The image of \(\widehat N\) in the original loop
abelianization is the lattice
\[
 \begin{aligned}
 \Lambda
   &=\{(a,b)\in\widehat{\mathbb Z}^2:
           a\equiv b\equiv0\pmod2,\ a+b\equiv0\pmod3\}\\
   &=\widehat{\mathbb Z}(2,-2)
                           +\widehat{\mathbb Z}(0,6).
 \end{aligned}                                                   \tag{D.47.5}
\]
Thus \(\widehat N/M=\Lambda\), and, for the actual
finite base \(H_m=\widehat B/L_m\),
\[
 \begin{aligned}
 c_m&=\gcd(2,m)^2\gcd(3,m),\\
 |H_m|&=576\,m^2/c_m,\qquad
 \operatorname{rank}(L_m)=1+96\,m^2/c_m .
 \end{aligned}                                                   \tag{D.47.6}
\]
The rank is that of its original finite-index discrete
free subgroup, whose completion is \(L_m\).
In particular \(H_{12}\) and \(L_{12}\) are exactly
the order-\(6912\), rank-\(1153\) source of D.37--D.39.

The actual Galois commutator \(h\) of D.46 fixes
\(\mathcal H_\infty\) pointwise. There is a continuous
peripheral automorphism of \(\mathcal V\)
\[
             \widetilde\alpha^\dagger
                =(\alpha_{h,\infty,2},
                            (1_{\mathcal P_{\infty,p}})_{p\ne2}),
                                                               \tag{D.47.7}
\]
with multiplier one and one compatible pair of normalized
loop conjugators. Its parameter is the literal tuple
\[
          \widetilde f^\dagger=(f_{h,\infty,2},1,1,\ldots)
                            \in\operatorname{im}(M\to\mathcal V).
                                                               \tag{D.47.8}
\]
It projects to the entire D.46 action, including the
prescribed finite pair \(T_4=I,\ W_3=[i]\).
It induces compatible actions on every Fermat homology
group at every deck level and every coefficient prime:
the full two-primary component is the actual \(h\)
action, and every odd-primary component is identity.
The same assertion holds for the congruence/Fermat
subgroups \(L_m\). It consequently supplies a replacement
on the same underlying all-prime homology source
\(\mathcal T\) used in D.39, with its new arithmetic
\(T_3\) component; the excluded D.39 action is not
retained.

In particular the exact inclusions
\[
                    R_\infty\subset M'
                                     \subset\widehat F''      \tag{D.47.9}
\]
give a quotient map onto the full metabelian group
\(\widehat F/\widehat F''\), and the action descends
under this map. Write
\[
 A=\widehat F'/\widehat F''=\prod_pA_p,\qquad
 \mathcal B_{\mathrm{met},p}
       =\widehat B/\ker(\widehat F'\to A_p).
\]
The exact metabelian action and parameter are
\[
 \begin{aligned}
 \widehat B/\widehat F''
       &=\prod_{\widehat B/\widehat F'}
                                      \mathcal B_{\mathrm{met},p},\\
 \alpha_{\mathrm{met}}^\dagger
       &=(\alpha_{h,\mathrm{met},2},
                                (1)_{p\ne2}),\\
 f_{\mathrm{met}}^\dagger
       &=(f_{h,\mathrm{met},2},(1)_{p\ne2}).
 \end{aligned}                                                   \tag{D.47.10}
\]
Here the prime is the coefficient prime; the complete
profinite deck group is retained in every component.
This metabelian action is nontrivial and its outer
class is not in the actual Galois image. In the
multiplier-one normalization used here, passage to
the outer action is injective. Under the two marked cover
maps of D.36 its compact coordinates are
\[
       \rho_{C_3,2}=\eta_*,\qquad
       \rho_{C_3,4}=-I,\qquad
       \rho_{C_4,3}=I,                                        \tag{D.47.11}
\]
which violate the proved arithmetic joint graph.
This is a proved comparison for the full metabelian
quotient and the specified relative source. Whether
this metabelian action, or its prescribed CM pair,
is induced by a genuine full profinite GT0 element,
or this metabelian action by a full loop-preserving
automorphism of \(\widehat F\), remains undecided.

The maps back to D.46 also have exact kernels.
Let \(\Lambda_p\) be the \(p\)-component of (D.47.5), and put
\[
 C_p=\ker(N^{(p)}\to\Lambda_p),\qquad
 D_p=\ker(M^{(p)}\twoheadrightarrow C_p).
                                                               \tag{D.47.12}
\]
These maps are the original subgroup inclusions
followed by the loop abelianization. Then
\[
 \begin{gathered}
 \ker(\mathcal U\to\mathcal H_\infty)=\prod_pC_p,\qquad
 1\longrightarrow\prod_pD_p\longrightarrow\mathcal V
                    \longrightarrow\mathcal U\longrightarrow1,\\
 \ker(\mathcal V\to\mathcal P_{\infty,p})
                         =\prod_{q\ne p}M^{(q)},\qquad
 \ker(\widehat B\to\mathcal V)=R_\infty .
 \end{gathered}                                               \tag{D.47.13}
\]
Under the actual \(S_7\) source (D.46.9), \(R_\infty\)
still has exact image \(A_7\). Thus the full-source
lifting step has not been replaced by a cofinality claim.

These results continue Part D(e), section 51.16,
folio 734 verso, scan 174, and the LM2 two-torus
and Fermat-tower follow-ups, section 51.15,
folios 725--727, scans 156--160.
\end{statement}
\begin{proof}
\emph{The marked base and its integral lattice.}

The program \texttt{d\_cm\_fermat\_relative\_tower.py}
enumerates \(Q\) from its original \(s,r\) matrices.
For \(x=(sr)^{-2}\), \(y=sxs\), it obtains
\(|H|=96\) and \(|H'|=8\).
It evaluates the actual conjugates of
\([x,y]\), checks the entire derived image,
and verifies the quaternion multiplication
conditions. The induced abelian quotient
has type \(C_2\times C_6\).
For all \(144\) pairs of exponents modulo
twelve the certificate proves
\[
          x^a y^b\in H'
        \quad\Longleftrightarrow\quad
          2\mid a,\quad2\mid b,\quad3\mid(a+b).
                                                               \tag{D.47.14}
\]
Right exactness of abelianization identifies
the image of \(N\) in \(F^{\rm ab}\)
with precisely this kernel lattice.
Its columns \((2,-2),(0,6)\) have
determinant twelve and Smith invariants
\(2,6\). Completion gives (D.47.5)
with its actual source map.

The source kernels of
\(\widehat B\to Q\) and
\(\widehat B\to\widehat B/\widehat F'\)
have product \(\widehat N\widehat F'\);
their maximal common marked quotient
is \(Q/q(\widehat F')=Q/J\).
Their intersection is \(M\).
Closed Goursat proves both directions
of (D.47.4). The image of
\(L_m\) under \(\widehat N\to\Lambda\)
is \(\Lambda\cap m\widehat{\mathbb Z}^2\).
Its quotient has order
\[
 |\Lambda/(\Lambda\cap m\widehat{\mathbb Z}^2)|
       =m^2/\bigl(\gcd(2,m)\gcd(6,m)\bigr)=m^2/c_m .
\]
This proves the base orders in
(D.47.6). These subgroups are
torsion-free finite-index subgroups
of \(C_2*C_3\); their actual
Bass--Serre graph gives
rank \(1+[B:L_m]/6\).
The certificate independently
enumerates the lattice modulo
every \(1\le m\le60\), including
the \(m=4,12\) comparisons.

\emph{Both inclusions in the closed-source inverse limit.}
Let \(R_{m,p}=\ker(L_m\to L_m^{(p)})\).
The topology and quotients here
are continuous profinite ones.
The actual restriction
\(M\to L_m^{(p)}\) has pro-\(p\)
image, so
\[
                    R_{\infty,p}\subset R_{m,p}
                              \quad\hbox{for every }m.        \tag{D.47.15}
\]
Conversely let \(C\) be the kernel
of any finite continuous
\(p\)-group quotient of \(M\).
Conjugation by \(\widehat B\)
has a finite orbit on \(C\).
Indeed \(C\) is clopen in
the compact group \(M\);
uniform continuity of
conjugation gives an open
neighbourhood of one in
\(\widehat B\) that preserves
its finite coset partition.
The normal core formed by
the finitely many conjugates
is still open in \(M\),
normal in \(\widehat B\),
and has \(p\)-group quotient.
Replace \(C\) by that core.

Inside the profinite group
\(\widehat B/C\), the subgroup
\(M/C\) is finite. There is
an open normal subgroup
disjoint from its nonidentity
elements. Its preimage \(O\)
satisfies
\[
                       O\cap M=C .
\]
In the finite quotient
\(\widehat B/O\), the image
of \(M\) is a \(p\)-group.
The quotient by that image
is a finite quotient of
\(\mathcal H_\infty\).
The subgroups \(L_m/M\)
are cofinal among its open
normal source kernels:
they are directed under
least common multiples,
and their intersection
is one. Compactness
therefore gives an \(m\)
with \(L_m\subset OM\).
The image of \(L_m\)
in \(\widehat B/O\)
is a \(p\)-group.
It follows that
\(R_{m,p}\subset O\).

An element of
\(\bigcap_mR_{m,p}\)
already lies in
\(\bigcap_mL_m=M\);
it consequently lies
in \(O\cap M=C\).
Applying this to every
finite \(p\)-quotient
of \(M\) proves
\[
                  \bigcap_mR_{m,p}
                          =R_{\infty,p}.                     \tag{D.47.16}
\]
For \(m\mid n\),
the restriction to
\(L_n\) of the
pro-\(p\) quotient
of \(L_m\) gives
\(R_{n,p}\subset R_{m,p}\).
The inverse system is
therefore an actual
system of quotients
of \(\widehat B\).
Compactness and
(D.47.16) prove the
first line of
(D.47.3), not just
one of its projections.

\emph{The maximal common source and all projection kernels.}
The continuous image of
the compact group \(M\)
in a product of finitely
many finite \(p\)-group
quotients at distinct
primes is the full
product. Goursat gives
this inductively from
coprime orders.
Compactness over all
levels proves
\[
                 M/R_\infty\simeq\prod_pM^{(p)} .
                                                               \tag{D.47.17}
\]
For two distinct
primes, the same
argument gives
\(R_{\infty,p}
 R_{\infty,q}=M\).
For one prime and
all remaining primes
it gives the analogous
product-kernel equality.
These are exactly
the two maximal-common
quotient assertions
over \(\mathcal H_\infty\).
The marked fibre
product in (D.47.3)
follows from the
same original source
map and (D.47.17).

The D.46 group
\(\mathcal U\) retains
all of
\(\widehat F^{\rm ab}\),
so its source kernel
lies in \(M\), and
it maps onto
\(\mathcal H_\infty\).
The image of \(M\)
in \(N^{(p)}\) is
exactly \(C_p\).
To see both directions,
its quotient is the
maximal pro-\(p\)
quotient of
\(\widehat N/M=\Lambda\),
which is \(\Lambda_p\).
The compact image of
\(M\) is closed, so
this identifies its
kernel without a
density assumption.
The induced
\(M^{(p)}\to C_p\)
is onto. The same
restriction
\(M\to N^{(p)}\)
also shows
\(R_{\infty,p}\subset R_p\);
hence
\(R_\infty\subset R_{\rm nil}\)
and gives the actual
source map
\(\mathcal V\to\mathcal U\).
Applying
the distinct-prime
product argument
again identifies
the image of \(M\)
in \(\mathcal U\)
with \(\prod_pC_p\).
The source map
\(\mathcal V\to\mathcal U\)
is the identity on
\(\mathcal H_\infty\)
and has kernel
\(\prod_pD_p\).
This proves every
map and kernel in
(D.47.12)--(D.47.13).

\emph{One coherent pair of normalized conjugators.}
For the actual \(h\)
of D.46, the normalized
action on \(Q\) is
the identity. Its
full multiplier is
one and \(f_h\in
\widehat F'\), so
the exact equations
\(\alpha_h(s)=s f_h\),
\(\alpha_h(r)=f_h^{-1}r\)
make its action on
\(\widehat B/\widehat F'\)
the identity as well.
Both coordinates in
(D.47.4) are fixed;
hence its action on
\(\mathcal H_\infty\)
is exactly the identity.
It preserves \(M\)
and every characteristic
pro-\(p\) residual,
giving the arithmetic
factor in (D.47.7).
The product action
is thus an actual
continuous automorphism
of (D.47.3).

Let \(A_\infty,\beta_\infty\)
be the images in
\(\mathcal H_\infty\)
of the normalized
Galois conjugators
of \(r,s\).
They centralize
the respective
elements. There is
even an exact
finite identification
of \(A_\infty\).
On the loop
abelianization,
the matrix of \(r\)
has
\(\det(1-r)=3\);
the program computes
that matrix from
the actual sphere
graph. The operator
\(1-r\) is injective
on
\(\widehat{\mathbb Z}^2\).
Since \(A_\infty\)
centralizes \(r\),
its abelian coordinate
is zero. It therefore
lies in \(q(\widehat F')
=Q_8\). The certificate
computes
\[
 C_{Q_8}(r)=\{I,5I\}.
\]
The arithmetic sign
\(T_{h,4}=-I\)
selects \(A_\infty=5I\)
under the exact
kernel identification
in (D.47.4).
In particular this
is a nontrivial
central element of
the marked base.

Apply the pro-two
coprime order-three
coset lifting proof
of D.46 to the map
\(\mathcal P_{\infty,2}
\to\mathcal H_\infty\).
There is one element
centralizing \(r\)
above \(A_\infty^{-1}\).
Right multiplication
of the arithmetic
\(r\)-conjugator
by it leaves the
automorphism unchanged
and makes its entire
\(\mathcal H_\infty\)
coordinate one.
At every odd prime,
the same coprime
order-two proof
gives an actual
\(s\)-centralizing
element above
\(\beta_\infty\).
These are lifts
of the whole
profinite base
elements, not
separate incompatible
finite choices.
The proof by finite
complements and
compactness works
over this infinite
base because its
kernel is pro-\(p\).

The resulting tuples
of conjugators belong
to the actual fibre
product (D.47.3).
They give (D.47.7),
fix \(t\), and give
exactly (D.47.8).
The parameter has
trivial base image,
so every source
preimage belongs
to \(M\subset
\widehat F'\).
The two hexagons
and generation
follow as in D.46.
Their projections
to \(\mathcal U\)
give the same
automorphism and
parameter as D.46,
even when a different
centralizer representative
was chosen there.
The invariant
\(G_2/G_3\)
framing map therefore
retains \(T_4=I\),
\(W_3=[i]\).
This constructs
one simultaneous
lift at every
prime and every
Fermat level.

\emph{Every homology quotient and the full metabelian action.}
An element killed by
every pro-\(p\)
quotient of \(M\)
is killed by every
finite abelian quotient
of \(M\): a finite
abelian group is a
product of its
prime-primary factors.
It therefore lies
in the closed
commutator \(M'\).
Since \(M\subset
\widehat F'\),
we obtain the
exact inclusions
(D.47.9).
Also \(M\subset L_m
\subset F_m\), so
\[
                R_\infty\subset M'\subset L_m'
                                               \subset F_m'. \tag{D.47.18}
\]
This proves every
claimed quotient
map from the source.

For action descent,
the quotient
\(\mathcal H_\infty\)
retains the entire
loop abelianization,
and (D.47.7) fixes
that base. Hence
the images of
\(F_m,L_m\) are
invariant; so are
their closed
commutators.
Their homology
quotients are actual
automorphism quotients.
For a specified
coefficient prime
\(p\), the restriction
of each such quotient
map to \(M\) has
abelian pro-\(p\)
image after retaining
the relevant finite
deck base. It thus
kills
\(R_{\infty,p}\).
The action on
that coefficient
component is exactly
the \(h\) action
at two and identity
at every odd prime.
The universal
property and the
finite-level maps
in (D.47.3) make
these statements
compatible for all
deck and coefficient
levels.

The subgroup
\(\widehat F\)
is the preimage
of the fixed
\(S_3\) quotient.
Its image in
\(\mathcal V\)
is invariant.
Its first and
second closed
commutators are
therefore invariant
as well, so
(D.47.9) gives
action descent to
the full metabelian
quotient.
The map from \(M\)
to each \(A_p\)
is pro-\(p\);
hence
\(\mathcal B_{\mathrm{met},p}\)
is a quotient of
\(\mathcal P_{\infty,p}\).
Decomposing the
full abelian
kernel
\(\widehat F'/\widehat F''\)
into its prime
factors proves
the first line of
(D.47.10) by its
actual source map.
The other lines
now follow from
(D.47.7)--(D.47.8).
They give the
complete metabelian
action, including
all prime-to-\(p\)
deck coordinates,
not merely the
pure pro-\(p\)
quotient of \(F\).

For the earlier
D.39 underlying
source, (D.47.18)
at \(m=12\)
gives its full
\(L_{12}\)-homology
factor. The map
to the actual
relative pro-two
source gives
its \(T_3\) factor.
Both are equivariant
under the replacement
action; their joint
source image is
the same underlying
\(\mathcal T\) by
D.39.2. The
two-primary homology
and \(T_3\) components
are now induced by
the actual \(h\);
the odd homology
components are
identity. This
specifies the
replacement under
the exact maps
without identifying
it with the
excluded earlier
action.

\emph{The actual CM information and nonarithmetical comparison.}
For the cyclic
three and four
subgroups \(K_3,K_4\)
of D.36, the
containment
\(\widehat F'\subset
\widehat K_n\)
gives
\(\widehat F''\subset
\widehat K_n'\).
Restriction to those
cover homology groups
and filling their
actual cusp lattices
therefore factor
through the full
metabelian source.
Their compact
two- and three-adic
modules factor
through
\(\mathcal P_{\infty,2}\)
and
\(\mathcal P_{\infty,3}\),
respectively, by
the just-proved
coefficient-prime
argument.
The exact maps,
cubic twist and
quartic tangent
convention are
those of D.34
and D.36. They
give (D.47.11).
In particular the
metabelian action
is nonidentity.

A Galois element
inducing this
metabelian action
would have full
cyclotomic character
one, from its
action on the
retained loop
abelianization.
It would lie
over
\(L=\mathbb Q(\zeta_{12})\).
The cubic matrix
modulo four has
trivial order-three
factor and
normalized sign
\(-1\).
D.36's exact
arithmetic graph
then forces
projective quartic
three-torsion
action \([i]\).
The actual
quartic coordinate
in (D.47.11)
is \(I\).
This contradiction
proves that the
metabelian action
is not in the
Galois image.
It proves the
comparison through
specific cover
maps and the
shared torsion
field, and does
not impose an
unproved full-GT0
cover-to-fixed-fibre
identity.

\emph{The outer comparison and normalization rigidity.}
Both the displayed
action and a
normalized Galois
action with multiplier
one fix \(x\) and
send \(y\) to
\(f^{-1}yf\), with
\(f\in\widehat F'\).
They are therefore
identity on the
free nilpotent
quotient of class
two. An inner
conjugator between
two such actions
acts trivially on
that quotient.
Its two abelian
coordinates vanish:
in each actual
Heisenberg quotient
\(U_3(\mathbb Z/n)\),
commuting with the
images of both
\(x\) and \(y\)
forces both coordinates
to be zero modulo
\(n\). Thus the
conjugator belongs
to
\(A=\widehat F'/\widehat F''
\simeq R_Ac\), where
\(R_A=\widehat{\mathbb Z}
[[\widehat{\mathbb Z}^{\,2}]]\),
by B.8. Fixing
\(x\) gives
\((X-1)a=0\)
in this module.

Here multiplication
by \(X-1\) is
injective in the
full profinite
completed group ring,
including every
prime-to-coefficient
deck component.
For a direct proof,
project to coefficients
\(\mathbb Z/m\) and
deck group
\((\mathbb Z/n)^2\),
and lift to deck
level \(mn\).
An \(X\)-invariant
element at that
higher level is
a sum of the
norm
\(1+X+\cdots+X^{mn-1}\)
times elements in
the other deck
coordinate. Its
pushforward to
level \(n\) is
multiplied by \(m\),
so is zero modulo
\(m\). This holds
for every \(m,n\),
proving injectivity.
The inner conjugator
is consequently
identity. The
normalized multiplier-one
actions embed in
the outer action
space. Since an
outer Galois comparison
already has multiplier
one on the retained
abelianization, it
would give the
based comparison
excluded above.
This proves the
claimed nonarithmetical
outer class as well.

Finally
\(\mathcal H_\infty\)
is prosolvable,
by (D.47.4),
and the kernel
in (D.47.17)
is pronilpotent.
Every finite
quotient of
\(\mathcal V\)
is solvable.
Under the actual
\(S_7\) source
of D.46, the
quotient by the
image of
\(R_\infty\)
is consequently
solvable.
That image contains
\(A_7\); it is
contained in
\(A_7\) because
\(R_\infty\subset
\widehat N\subset
B_{\rm even}\).
It is exactly
\(A_7\), completing
the remaining
source-kernel
comparison.
\end{proof}
